% Source assembly: Mumford, Abelian Varieties, Chapter II Sections 7--9,
% source pages 77--98 (with the Section 6 applications on pages 72--76).

\section{Applications of the theorem of the square}
\label{mumford-source-II6-applications}
\dcref{M:ch2:II.6}
\uses{mumford-thm-seesaw,mumford-thm-cube-varieties}

For a line bundle $L$ on an abelian variety $X$, put
$\phi_L(x)=[T_x^*L\otimes L^{-1}]$ and
$K(L)=\ker\phi_L=\{x:T_x^*L\simeq L\}$.
\source{mumford-abelian-varieties:page-0071}

\begin{proposition}[Closedness of the translation kernel]
\label{mumford-prop-KL-closed}
\dcref{M:ch2:II.6}
\uses{mumford-thm-seesaw}
$K(L)$ is a Zariski-closed subgroup of $X$.
\source{mumford-abelian-varieties:page-0071}
\end{proposition}
\begin{proof}
Apply the seesaw theorem to $m^*L\otimes p_2^*L^{-1}$ on $X\times X$.
The locus of $x$ for which its restriction to $\{x\}\times X$ is trivial is
closed, and that restriction is $T_x^*L\otimes L^{-1}$.  It is therefore
exactly $K(L)$.
\end{proof}

\begin{proposition}[Effective-divisor criterion]
\label{mumford-prop-effective-divisor-criterion}
\dcref{M:ch2:II.6}
\uses{mumford-prop-KL-closed,mumford-thm-seesaw}
Let $D$ be an effective divisor on $X$, $L=\mathcal O_X(D)$, and
$H=\{x:T_x^*D=D\}$.  The following are equivalent:
\begin{enumerate}
\item $H$ is finite;
\item $K(L)$ is finite;
\item $|2D|$ is base-point free and defines a finite morphism
      $X\to\mathbf P^N$;
\item $L$ is ample.
\end{enumerate}
\source{mumford-abelian-varieties:page-0071}\source{mumford-abelian-varieties:page-0072}\source{mumford-abelian-varieties:page-0073}
\end{proposition}
\begin{proof}
The implications (iii)$\Rightarrow$(iv) and (ii)$\Rightarrow$(i) are the
standard finite-morphism criterion and $K(L)\supset H$.  If $K(L)$ is not
finite, let $Y=K(L)^0$.  The restriction $L_Y$ is ample and invariant under
all translations by $Y$.  The seesaw theorem makes
$m^*L_Y\otimes p_1^*L_Y^{-1}\otimes p_2^*L_Y^{-1}$ trivial; pullback by
$y\mapsto(y,-y)$ makes $L_Y\otimes(-1_Y)^*L_Y$ trivial.  This is impossible
because it is ample on the positive-dimensional $Y$, proving (iv)$\Rightarrow$(ii).

Assume (i).  The divisors $T_x^*D+T_{-x}^*D$ lie in $|2D|$.  Given $v\in X$,
choose $x$ with $v\pm x\notin\operatorname{Supp}D$; hence $|2D|$ has no
base points and defines $\phi:X\to\mathbf P^N$.  If $\phi$ were not finite,
there would be an irreducible curve $C$ contracted by $\phi$, so every member
of $|2D|$ would contain $C$ or be disjoint from it.  For almost all $x$,
$C$ is disjoint from $T_x^*D+T_{-x}^*D$.

We use the following lemma.  If $C$ is a curve and $E$ an irreducible divisor
with $C\cap E=\varnothing$, then $E$ is invariant under translation by
$x_1-x_2$ for all $x_1,x_2\in C$.  Indeed, $\mathcal O_X(E)|_C$ is trivial,
so every translate has degree zero on $C$.  Thus a translate of $C$ either
misses $E$ or is contained in $E$.  For $y\in E$, the translate
$T_{y-x_2}(C)$ meets $E$ at $y$ and is consequently contained in $E$;
therefore $y-x_2+x_1\in E$.

Writing $D=\sum D_i$, the lemma makes every $D_i$ invariant under all
$x_1-x_2$ with $x_i\in C$, contradicting (i).  Thus $\phi$ is finite and
(i)$\Rightarrow$(iii).
\end{proof}

\begin{theorem}[Projectivity of abelian varieties]
\label{mumford-thm-abelian-projective}
\dcref{M:ch2:II.6}
\uses{mumford-prop-effective-divisor-criterion}
Every abelian variety is projective.
\source{mumford-abelian-varieties:page-0073}
\end{theorem}
\begin{proof}
Choose an affine open $U\subset X$ and let $D_1,\ldots,D_r$ be the
irreducible components of $X-U$, with $D=\sum D_i$.  Translate so that
$0\in U$.  The subgroup $H=\{x:T_x^*D=D\}$ is closed, and each $x\in H$
preserves $U$, hence $H\subset U$.  Since $H$ is complete and $U$ is affine,
$H$ is finite.  The criterion therefore makes $\mathcal O_X(D)$ ample, so a
power embeds $X$ in projective space.
\end{proof}

\begin{proposition}[Hilbert-polynomial degree]
\label{mumford-prop-degree-hilbert}
\dcref{M:ch2:II.6}
Let $X$ be complete of dimension $g$, $L$ invertible, and $\mathcal F$ coherent.
Write $P_{\mathcal F}(n)=\chi(\mathcal F\otimes L^n)$ and
$d_L(\mathcal F)/g!$ for its leading coefficient.  Then
\[
d_L(\mathcal F)=\operatorname{rank}(\mathcal F)\deg L.
\]
If $f:Y\to X$ is a surjective morphism of complete $g$-folds, then
\[
\deg f^*L=(\deg f)\deg L.
\]
\source{mumford-abelian-varieties:page-0075}\source{mumford-abelian-varieties:page-0076}
\end{proposition}
\begin{proof}
Choose a nonzero ideal sheaf $\mathcal J$ with an exact sequence
$0\to\mathcal J^{\operatorname{rank}\mathcal F}\to\mathcal F\to\mathcal G\to0$
where $\mathcal G$ is torsion.  Its support has dimension $<g$, so it does
not affect the leading coefficient.  Additivity of $\chi$ and
$0\to\mathcal J\to\mathcal O_X\to\mathcal O_X/\mathcal J\to0$ give
$d_L(\mathcal F)=\operatorname{rank}(\mathcal F)d_L(\mathcal J)
=\operatorname{rank}(\mathcal F)\deg L$.

For (2), projection formula gives
$R^pf_*(f^*L^n)=R^pf_*\mathcal O_Y\otimes L^n$.  The Leray spectral sequence
therefore yields
\[
\chi(f^*L^n)=\sum_p(-1)^p\chi(R^pf_*\mathcal O_Y\otimes L^n).
\]
For $p>0$ the higher direct images are supported on a proper closed subset,
while $f_*\mathcal O_Y$ has rank $\deg f$.  Comparing the coefficients of
$n^g$ proves the formula.
\end{proof}

\begin{corollary}[Degree of multiplication]
\label{mumford-cor-degree-multiplication-appendix}
\dcref{M:ch2:II.6}
\uses{mumford-prop-degree-hilbert,mumford-thm-cube-varieties}
For an abelian variety of dimension $g$, $\deg(n_X)=n^{2g}$.
\source{mumford-abelian-varieties:page-0076}
\end{corollary}
\begin{proof}
Choose an ample symmetric line bundle $L$.  The theorem of the square gives
$n_X^*L\simeq L^{n^2}$, hence
$\deg n_X\deg L=\deg(n_X^*L)=\deg(L^{n^2})=n^{2g}\deg L$.
Since $\deg L>0$, cancel to obtain the assertion.
\end{proof}

\section{Dividing varieties by finite groups}
\label{mumford-source-II7}
\dcref{M:ch2:II.7}
\uses{mumford-def-abelian-variety}

\begin{theorem}[Finite group quotient]
\label{mumford-thm-finite-group-quotient}
\dcref{M:ch2:II.7}
\uses{mumford-def-abelian-variety}
Let a finite group $G$ act on a variety $X$, with every orbit contained in an
affine open. Then there is a quotient variety $Y=X/G$ with
$\mathcal O_Y=\pi_*\mathcal O_X^G$; $\pi$ is finite, surjective, separable,
and etale when the action is free. If $X$ is affine, so is $Y$.
\source{mumford-abelian-varieties:page-0077}\source{mumford-abelian-varieties:page-0078}\source{mumford-abelian-varieties:page-0079}
\end{theorem}
\begin{proof}
If $U'$ is affine and contains an orbit $Gx$, then
$U=\bigcap_{g\in G}gU'$ is affine, $G$-stable, and contains $x$.  Thus it
suffices to prove the assertion for $X=\operatorname{Spec}A$.  Write
$A=k[x_1,\ldots,x_n]$, $v=|G|$, and let $B=A^G$.  If $\sigma_j(x_i)$ is
the $j$th elementary symmetric function in the orbit of $x_i$, put
$B'=k[\sigma_j(x_i)]$.  Each $x_i$ satisfies
\[
 T^v-\sigma_1(x_i)T^{v-1}+\cdots+(-1)^v\sigma_v(x_i)=0,
\]
so $A$ is finite over the noetherian ring $B'$.  Since
$B'\subset B\subset A$, $B$ is finite over $B'$ and $A$ is finite over $B$.
The inclusion $B\hookrightarrow A$ therefore defines a finite surjective map
$\pi:\operatorname{Spec}A\to\operatorname{Spec}B$.

If $R=K(A)$, then $R^G=K(B)$: for $a/b\in R^G$,
\[
 \frac ab=a\left(\prod_{g\ne e}g(b)\right)\left(\prod_{g\in G}g(b)\right)^{-1},
\]
and the numerator and denominator on the right are invariant.  Thus the
extension of quotient fields is separable.  The invariant direct image is the
intersection of the kernels of $f\mapsto g(f)-f$, and the affine calculation
identifies it with $\mathcal O_Y$.

If two points have distinct $G$-orbits, choose $f\in A$ equal to $1$ on one
orbit and $0$ on the other.  The product $\prod_{g\in G}g(f)$ is invariant and
separates their images, so the finite map has the orbit quotient topology.

Finally, complete at $y=\pi(x)$.  The Chinese remainder theorem gives
\[
 \widehat A\simeq\prod_{g\in G}\widehat{\mathcal O}_{X,gx}.
\]
The exact sequence
\[
0\to B\to A\to\prod_{g\in G}A,
\qquad a\mapsto(g(a)-a)_g,
\]
remains exact after the flat completion $\widehat B$.  Identifying all
completed local factors by the action, $G$ permutes the factors, and its
invariants are the diagonal copy of $\widehat{\mathcal O}_{X,x}$.  Hence
$\widehat{\mathcal O}_{Y,y}\to\widehat{\mathcal O}_{X,x}$ is an isomorphism
when the action is free, proving that $\pi$ is etale.  The affine quotients
glue over the $G$-stable affine cover.
\end{proof}

\begin{proposition}[Descent of coherent sheaves]
\label{mumford-prop-coherent-descent-finite-quotient}
\dcref{M:ch2:II.7}
\uses{mumford-thm-finite-group-quotient}
If $G$ acts freely on $X$ and $Y=X/G$, pullback is an equivalence between
coherent sheaves on $Y$ and coherent $G$-sheaves on $X$.  Its inverse is
$\mathcal G\mapsto\pi_*\mathcal G^G$, and locally free sheaves correspond with
the same rank.
\source{mumford-abelian-varieties:page-0080}\source{mumford-abelian-varieties:page-0081}\source{mumford-abelian-varieties:page-0082}
\end{proposition}
\begin{proof}
There are natural maps
\[
 S(\mathcal F):\mathcal F\to\pi_*(\pi^*\mathcal F)^G,qquad
 T(\mathcal G):\pi^*(\pi_*\mathcal G^G)\to\mathcal G.
\]
Affinely, with $X=\operatorname{Spec}A$, $Y=\operatorname{Spec}B$,
$B=A^G$, these become
\[
 S(M):M\to(M\otimes_BA)^G,qquad
 T(N):(N^G)\otimes_BA\to N.
\]
The composite after tensoring $S(M)$ with $A$ is the identity.  Since $A$ is
faithfully flat over $B$, it is enough to prove that all $T(N)$ are isomorphisms.
After completion at a point of $Y$, freeness identifies
$A\otimes_B\widehat B$ with a product of copies of $\widehat B$, one for each
element of $G$, with $G$ acting by simply transitive permutation.  In this
model the map $T(N)$ is immediate; faithful flatness of completion descends
the isomorphism.  The same argument preserves finite projective modules and
therefore local freeness and rank.
\end{proof}

\begin{proposition}[Translation quotients]
\label{mumford-prop-translation-quotient-isogeny}
\dcref{M:ch2:II.7}
\uses{mumford-thm-finite-group-quotient}
For a finite subgroup $K$ of an abelian variety $X$, the quotient $X/K$ is an
abelian variety and $X\to X/K$ is an etale isogeny. Conversely, every isogeny
is a quotient by its finite kernel.
\source{mumford-abelian-varieties:page-0083}\source{mumford-abelian-varieties:page-0084}\source{mumford-abelian-varieties:page-0085}
\end{proposition}
\begin{proof}
Translations by $K$ give a free finite-group action, so the quotient theorem
gives a finite etale map $X\to X/K$. The group law on the abstract quotient
is a morphism by applying the quotient property to the addition map; inversion
is treated similarly. Conversely, an isogeny is the quotient by its finite
kernel.
\end{proof}

\begin{proposition}[Characters and descended line bundles]
\label{mumford-source-character-kernel-pic}
\dcref{M:ch2:II.7}
\uses{mumford-thm-finite-group-quotient}
If $G$ acts freely on $X$ and $Y=X/G$, then for characters $\alpha:G\to k^*$,
\[
 \widehat G\simeq\ker[\Pic Y\to\Pic X],
\]
and if $|G|$ is prime to the characteristic, every coherent sheaf on $Y$ is
a direct summand of $\pi_*(\pi^*\mathcal F)$.
\source{mumford-abelian-varieties:page-0082}\source{mumford-abelian-varieties:page-0083}
\end{proposition}
\begin{proof}
An action on the trivial sheaf covering the action on $X$ sends the unit
section to a nowhere-vanishing regular function, hence to a scalar because
$X$ is complete. This scalar is a character; conversely every character
defines the action $g(f)=\alpha(g)^{-1}(f\circ g^{-1})$.
\end{proof}

\begin{corollary}[Divisibility and finite torsion]
\label{mumford-cor-divisible-finite-torsion}
\dcref{M:ch2:II.7}
An abelian variety $X$ is divisible and $X_n$ is finite for every $n>1$.
\source{mumford-abelian-varieties:page-0073}
\end{corollary}
\begin{proof}
For an ample line bundle $L$, the theorem of the square gives
\[
 n_X^*L\simeq L^{\otimes n(n+1)/2}\otimes(-1_X)^*L^{\otimes n(n-1)/2},
\]
which is ample. Since its restriction to $\ker(n_X)$ is trivial, the kernel
cannot contain a positive-dimensional subvariety. Thus it is finite, and the
surjectivity follows from the dimension theorem.
\end{proof}

\begin{proposition}[Degree and separability of multiplication]
\label{mumford-prop-multiplication-degree-separability}
\dcref{M:ch2:II.7}
For $g=\dim X$,
\[
 \deg(n_X)=n^{2g},\qquad n_X\text{ is separable}\Longleftrightarrow p\nmid n.
\]
If $p\nmid n$, then $X_n\simeq(\mathbf Z/n\mathbf Z)^{2g}$; for some
$0\le i\le g$, $X_{p^m}\simeq(\mathbf Z/p^m\mathbf Z)^i$.
\source{mumford-abelian-varieties:page-0073}\source{mumford-abelian-varieties:page-0074}\source{mumford-abelian-varieties:page-0075}
\end{proposition}
\begin{proof}
For equal-dimensional complete varieties,
$(f^*D_1\cdots f^*D_g)=\deg(f)(D_1\cdots D_g)$. For an ample symmetric
divisor $D$, $n_X^*D\sim n^2D$, hence $\deg(n_X)=n^{2g}$. If $p\nmid n$,
this degree is prime to $p$, so $n_X$ is separable and $\#X_n=n^{2g}$;
counting all $X_m$ for $m\mid n$ gives the stated group. If $p\mid n$,
$d(n_X)=0$, the rational differential map is zero, and the inseparable degree
is at least $p^g$. Thus $X_p\simeq(\mathbf Z/p\mathbf Z)^i$, and divisibility
gives the assertion for $p^m$.
\end{proof}

\section{The dual abelian variety in characteristic zero}
\label{mumford-source-II8}
\dcref{M:ch2:II.8}
\uses{mumford-thm-cube-varieties,mumford-thm-seesaw}

\begin{definition}[The dual variety]
\label{mumford-def-piczero}
\dcref{M:ch2:II.8}
\uses{mumford-thm-cube-varieties}
$\Pic^0(X)$ is the subgroup of line bundles $L$ for which
$\phi_L(x)=T_x^*L\otimes L^{-1}$ is identically zero. In characteristic zero
it has a natural structure of an abelian variety, denoted $\widehat X$.
\source{mumford-abelian-varieties:page-0085}\source{mumford-abelian-varieties:page-0086}
% Source introduces this definition and develops its structure in the following
% theorem; it gives no separate proof for the definition itself.
\notready
\end{definition}

\begin{theorem}[Dual abelian variety, characteristic zero]
\label{mumford-thm-dual-char-zero}
\dcref{M:ch2:II.8}
\uses{mumford-thm-cube-varieties,mumford-thm-seesaw,mumford-thm-finite-group-quotient}
For every abelian variety $X$ over a characteristic-zero field there is an
abelian variety $\widehat X$ representing $\Pic^0(X)$, with a Poincare bundle
$P_X$ on $X\times\widehat X$ trivial on both coordinate axes. Every line
bundle in $\Pic^0(X)$ is represented uniquely by a point of $\widehat X$.
\source{mumford-abelian-varieties:page-0086}\source{mumford-abelian-varieties:page-0087}\source{mumford-abelian-varieties:page-0088}\source{mumford-abelian-varieties:page-0089}\source{mumford-abelian-varieties:page-0090}\source{mumford-abelian-varieties:page-0091}
\end{theorem}
\begin{proof}
The theorem of the square gives, for $L\in\Pic^0(X)$,
\[
 T_x^*L\simeq L,\qquad (x+y)^*L\simeq T_x^*L\otimes T_y^*L\otimes L^{-1},
 \qquad n_X^*L\simeq L^n.
\]
Conversely, for arbitrary $L$ one has
\[
 n_X^*L\simeq L^{n^2}\otimes
 [L\otimes(-1_X)^*L^{-1}]^{(n-n^2)/2},
\]
and the bracketed factor lies in $\Pic^0(X)$ by the theorem of the square.
Any finite-order line bundle is therefore in $\Pic^0(X)$, since
$0=\phi_{L^n}(x)=\phi_L(nx)$ and $X$ is divisible.

If $L\in\Pic^0(X)$ is nontrivial, then $H^0(X,L)=0$: a nonzero section would
write $L=\mathcal O_X(D)$ with $D\ge0$, while
$L^{-1}\simeq(-1_X)^*L=\mathcal O_X((-1_X)^*D)$, forcing $D=0$.  If $k$ is
the least index with $H^k(X,L)\ne0$, pull back the isomorphism
$m^*L\simeq p_1^*L\otimes p_2^*L$ along $s_1(x)=(x,0)$.  The Kunneth
decomposition gives
\[
 H^k(X\times X,m^*L)\simeq
 \bigoplus_{i+j=k}H^i(X,L)\otimes H^j(X,L)=0,
\]
since in every summand one index is $<k$.  As $m\circ s_1=1_X$, the identity
on $H^k(X,L)$ factors through zero, a contradiction.  Thus all cohomology
of a nontrivial $L\in\Pic^0(X)$ vanishes.

Let $L$ be ample and put
\[
 K=m^*L\otimes p_1^*L^{-1}\otimes p_2^*(L^{-1}\otimes M^{-1})
\]
for $M\in\Pic^0(X)$.  The two Leray spectral sequences for the projections
$p_1,p_2:X\times X\to X$ have abutment $H^*(X\times X,K)$.  If
$M\simeq T_x^*L\otimes L^{-1}$, the restriction of $K$ to the $p_1$-fiber
over $x$ is a nontrivial member of $\Pic^0(X)$ and hence has no cohomology.
Therefore the first spectral sequence gives $H^*(X\times X,K)=0$.  For the
second, the support of every $R^kp_{2,*}K$ is contained in the finite set
$K(L)=\{x:T_x^*L\simeq L\}$.  The second spectral sequence then degenerates,
so all restrictions $K|_{X\times\{x\}}$ have zero cohomology.  At $x=0$ this
restriction is trivial and has nonzero $H^0$, a contradiction unless the
chosen $M$ is of the displayed form.  Hence $\phi_L:X\to\Pic^0(X)$ is
surjective.

Take $\widehat X=X/K(L)$, which exists by the finite-quotient theorem, and
let $\pi:X\to\widehat X$.  Put
\[
 K_0=m^*L\otimes p_1^*L^{-1}\otimes p_2^*L^{-1}.
\]
For $a\in K(L)$, translation by $(0,a)$ preserves $K_0$ because
$T_a^*L\simeq L$.  Normalize the resulting lift on
$\{0\}\times X\simeq L^{-1}(0)\times X$ to be
$(\lambda,x)\mapsto(\lambda,x+a)$.  The normalization is unique, and hence
the lifts form an action of $K(L)$.  Descent along $1_X\times\pi$ produces a
line bundle $P$ on $X\times\widehat X$ with
\[
 (1_X\times\pi)^*P\simeq K_0.
\]
Its fiber at $\pi(x)$ is $T_x^*L\otimes L^{-1}$, and its restriction to
$\{0\}\times\widehat X$ is trivial, so $P$ is normalized.

Finally, if $S$ is normal and $N$ is a line bundle on $X\times S$ whose
fibers lie in $\Pic^0(X)$ and whose restriction to $\{0\}\times S$ is trivial,
form
\[
 E=p_{12}^*N\otimes p_{13}^*P^{-1}
 \quad\text{on }X\times S\times\widehat X.
\]
The seesaw theorem makes the locus where
$E|_{X\times\{s,\alpha\}}$ is trivial a closed subset $\Gamma$ of
$S\times\widehat X$.  It is the graph of the set-theoretic classifying map
$f:S\to\widehat X$.  The projection $\Gamma\to S$ is bijective and, in
characteristic zero with $S$ normal, Zariski's Main Theorem makes it an
isomorphism.  Thus $f$ is a morphism and seesaw gives
$N\simeq(1_X\times f)^*P$.  This universal property identifies
$\widehat X$ with the representing abelian variety for $\Pic^0(X)$.
\end{proof}

\begin{proposition}[Biduality]
\label{mumford-prop-biduality-char-zero}
\dcref{M:ch2:II.8}
\uses{mumford-thm-dual-char-zero,mumford-source-character-kernel-pic}
There is a canonical isomorphism $X\simeq\widehat{\widehat X}$.
If $f:X\to Y$ is an isogeny, then $\widehat f:\widehat Y\to\widehat X$ is an
isogeny and $\ker f$ and $\ker\widehat f$ are canonically dual finite groups.
\source{mumford-abelian-varieties:page-0091}\source{mumford-abelian-varieties:page-0092}
\end{proposition}
\begin{proof}
The pullback map on Picard varieties sends $\Pic^0(Y)$ into $\Pic^0(X)$ and
defines $\widehat f$. If $f$ is an isogeny, the kernel of $\widehat f$ is the
kernel of pullback on Picard groups, which is dual to $\ker f$ by the finite
quotient character calculation. Equal dimensions then imply that $\widehat f$
is an isogeny. Applying this construction twice gives the canonical bidual
identification.
\end{proof}

\begin{definition}[Divisorial correspondence]
\label{mumford-source-divisorial-correspondence}
\dcref{M:ch2:II.8}
A divisorial correspondence between $X,Y$ is a line bundle on $X\times Y$
trivial on $\{0\}\times Y$ and $X\times\{0\}$.
\source{mumford-abelian-varieties:page-0092}
\end{definition}

\begin{proposition}[Characterization of the Poincare correspondence]
\label{mumford-source-divisorial-biduality}
\dcref{M:ch2:II.8}
For equal-dimensional $X,Y$ and a divisorial correspondence $Q$, the
conditions that $Q|_{\{x\}\times Y}$ be trivial only for $x=0$ and that
$Q|_{X\times\{y\}}$ be trivial only for $y=0$ are equivalent. When they hold,
$X\simeq\widehat Y$, $Y\simeq\widehat X$, and $Q$ is the Poincare bundle.
\source{mumford-abelian-varieties:page-0092}\source{mumford-abelian-varieties:page-0093}
\end{proposition}
\begin{proof}
By the universal property, (1) gives an injective $X\to\widehat Y$ and hence
an isomorphism in equal dimensions and characteristic zero. The induced map
$\psi:Y\to\widehat X$ must be injective: otherwise factor it through
$Y/K$ for a nonzero finite subgroup $K$. The correspondence then induces
$X\to\widehat{Y/K}$ whose composite with $\widehat{Y/K}\to\widehat Y$ is
an isomorphism, contradicting the nontrivial kernel of the latter dual
isogeny. Symmetry gives (2) from (1).
\end{proof}

\section{Analytic comparison over $\mathbf C$}
\label{mumford-source-II9}
\dcref{M:ch2:II.9}
\uses{mumford-thm-appell-humbert,mumford-def-piczero}

\begin{theorem}[Analytic translation formula]
\label{mumford-thm-analytic-dual-comparison}
\dcref{M:ch2:II.9}
\uses{mumford-thm-appell-humbert,mumford-def-piczero}
For $L(H,\alpha)$ on $X=V/U$,
\[
 T_{\pi(a)}^*L(H,\alpha)\simeq L(H,\alpha\gamma_a),
 \qquad \gamma_a(u)=e^{2\pi iE(a,u)}.
\]
Consequently
\[
 K(L)=U^\perp/U,
\quad U^\perp=\{a:E(a,u)\in\mathbf Z\ \forall u\in U\},
\]
and $L\in\Pic^0(X)$ iff $H=0$; $K(L)$ is finite iff $H$ is nondegenerate.
\source{mumford-abelian-varieties:page-0094}\source{mumford-abelian-varieties:page-0095}
\end{theorem}
\begin{proof}
Pull back the factor of automorphy under translation by $a$ and conjugate by
$e^{-\pi\overline{H(z,a)}}$. The multiplier changes by
$e^{\pi(H(a,u)-H(u,a))}=e^{2\pi iE(a,u)}$.
\end{proof}

\begin{proposition}[Analytic construction of the dual]
\label{mumford-source-analytic-dual}
\dcref{M:ch2:II.9}
For $X=V/U$, let
\[
 U'=\{l\in\overline T:\operatorname{Im}l(u)\in\mathbf Z\ \forall u\in U\}.
\]
Then $\widehat X=\overline T/U'$, and the Poincare bundle is
$L(H,\alpha)$ with
\[
 H((x_1,l_1),(x_2,l_2))=\overline{l_2(x_1)}+l_1(x_2),
 \qquad \alpha((u,l))=e^{-\pi i\operatorname{Im}l(u)}.
\]
The point represented by $l$ corresponds to $L(0,\alpha_l)$ with
$\alpha_l(u)=e^{2\pi i\operatorname{Im}l(u)}$.
\source{mumford-abelian-varieties:page-0095}\source{mumford-abelian-varieties:page-0096}\source{mumford-abelian-varieties:page-0097}
% Source presents this analytic construction directly, without a separate proof.
\notready
\end{proposition}

\begin{remark}[Normalization of the analytic dual identification]
\label{mumford-source-analytic-dual-normalization}
\dcref{M:ch2:II.9}
The two analytic descriptions of the map from $\overline T$ to
$\operatorname{Pic}^0(X)$ differ by the scalar normalization $2i$; explicitly,
the factor-of-automorphy description sends
\[
 l\longmapsto L(0,\alpha_l^*),
 \qquad \alpha_l^*(u)=e^{4\pi i\operatorname{Re}l(u)}.
\]
\source{mumford-abelian-varieties:page-0098}
\end{remark}
